\documentclass[11pt]{article}
\usepackage[margin=1in]{geometry}
\usepackage{amsmath, amssymb, amsthm, mathtools}
\usepackage{booktabs}
\usepackage{graphicx}
\usepackage{caption}
\usepackage[hidelinks]{hyperref}
\usepackage{microtype}

\newtheorem{theorem}{Theorem}
\newtheorem{proposition}{Proposition}
\newtheorem{lemma}{Lemma}
\theoremstyle{definition}
\newtheorem{definition}{Definition}
\newtheorem{remark}{Remark}

\newcommand{\E}{\mathbb{E}}
\newcommand{\R}{\mathbb{R}}
\newcommand{\CE}{\mathrm{CE}}
\newcommand{\DER}{\mathrm{DER}}
\newcommand{\DSR}{\mathrm{DSR}}
\newcommand{\Var}{\mathrm{Var}}
\newcommand{\KL}{\mathrm{KL}}
\newcommand{\CVaR}{\mathrm{CVaR}}
\newcommand{\sgn}{\operatorname{sgn}}

\title{\textbf{The Differential Entropic Reward: A Risk-Sensitive Objective\\
for Reinforcement-Learning Trading Agents}}
\author{Stefanos Drakos\\
\small AGEL AI I.K.E., Rhodes, Greece\\
\small ORCID: 0000-0001-7417-2444}
\date{\today}

\begin{document}
\maketitle

\begin{abstract}
\noindent
We introduce the \emph{Differential Entropic Reward} (DER), an online,
differentiable per-step reward for reinforcement-learning (RL) trading agents
derived from the entropic certainty equivalent of returns. A single risk-aversion
parameter $\theta$ controls the agent's sensitivity to \emph{all} cumulants of the
return distribution simultaneously --- mean, variance, skewness (crash risk), and
kurtosis (fat tails) --- in contrast to the Differential Sharpe Ratio (DSR), which is
structurally blind to moments beyond the second, and in contrast to
Conditional-Value-at-Risk objectives, which require auxiliary critics or sample
sorting. We prove that the DER recovers the risk-neutral mean as $\theta\to0$, reduces
to a mean--variance criterion at first order, admits an exact telescoping identity
to the realized entropic certainty equivalent, and possesses a Kullback--Leibler
robust (worst-case) dual; it is the discrete-time, model-free counterpart of the
risk-sensitive Hamilton--Jacobi--Bellman equation and uses precisely the exponential
risk criterion of the modern risk-sensitive RL regret theory. Each theoretical claim
is confirmed by an independent numerical check ($8/8$, seed $2025$). Empirically, on
a universe of $48$ real S\&P~500 equities, the DER reduces maximum drawdown relative
to the DSR for \emph{$100\%$} of assets (paired Wilcoxon $p\approx7\times10^{-15}$),
eliminating all drawdowns worse than $-20\%$, whereas the worst DSR drawdown reaches
$-50\%$. The reward is architecture-agnostic: transferred unchanged to actor--critic
and PPO, \emph{DER-PPO matches the risk control of CVaR-PPO (CPPO) at roughly a quarter
of the compute and without LLM signals}. We are deliberately explicit about scope: through a sequence of honest
ablations on a $470$-stock cross-sectional universe we show that flexible signal
learning overfits, that volume/order-flow features carry no measurable directional
information at the daily horizon, and that the value of the DER is risk \emph{shaping}
on top of a fixed signal, not alpha generation --- a more candid position than the
regime-dependent, drawdown-prone results reported for state-of-the-art LLM-augmented
risk-sensitive agents. All claims are reproduced by accompanying verification code.

\medskip
\noindent\textbf{Keywords:} reinforcement learning, algorithmic trading, risk-sensitive
control, entropic risk measure, reward design, drawdown.
\end{abstract}

\section{Introduction}\label{sec:intro}

Reinforcement learning (RL) has become a standard tool for sequential trading
decisions, spanning recurrent reinforcement learning (RRL)~\cite{moody2001}, deep
direct RL~\cite{deng2017}, value-based agents~\cite{theate2021}, portfolio
policies~\cite{jiang2017}, optimal execution~\cite{nevmyvaka2006}, and market
making~\cite{spooner2018}. A persistent difficulty, emphasized in the survey of
Hambly, Xu and Yang~\cite{hambly2023}, is \emph{risk}: agents optimizing raw or
Sharpe-like objectives tend to incur large, persistent drawdowns that make them
unsuitable for capital that must be protected. Even current LLM-augmented
risk-sensitive agents report regime-dependent performance and acknowledged
drawdowns~\cite{benhenda2025,dapo2025}.

\paragraph{What prior work has done.} Three threads define the landscape.
\emph{First}, direct-policy trading: Moody and Saffell~\cite{moody2001} introduced
recurrent reinforcement learning together with the Differential Sharpe Ratio, an
online reward that rewards the marginal change of a recursively estimated Sharpe
ratio; Deng et al.~\cite{deng2017} added deep feature learning, Th\'eate and
Ernst~\cite{theate2021} a deep $Q$-network agent, and Jiang et al.~\cite{jiang2017}
an end-to-end portfolio network, while task-specific agents addressed
execution~\cite{nevmyvaka2006} and market making~\cite{spooner2018}; the FinRL
ecosystem~\cite{finrl2020} consolidated these into standard environments and
baselines. \emph{Second}, risk-sensitive learning and control: the entropic risk
measure and its Kullback--Leibler dual are classical~\cite{follmer2002}, the
continuous-time risk-sensitive portfolio problem is well
developed~\cite{bielecki1999,flemingmce1995}, and a recent line establishes regret
guarantees for RL under exactly the exponential criterion we
adopt~\cite{fei2020,fei2021,noorani2022,borkar2002}; the competing applied functional
is Conditional Value-at-Risk~\cite{rockafellar2000}, used by current LLM-augmented
agents~\cite{benhenda2025,dapo2025} that nonetheless report regime-dependent returns
and unresolved drawdowns. \emph{Third}, the quantitative-strategy literature has
documented the genuine, if modest, sources of return that any trading agent must rely
on --- cross-sectional and time-series momentum~\cite{jegadeesh1993,moskowitz2012}
--- together with their characteristic failure mode, momentum
crashes~\cite{daniel2016}, and the volatility-scaling remedy that mitigates
them~\cite{moreira2017,lim2019}. Our work connects these threads: it brings the
exponential risk criterion of risk-sensitive RL into the direct-policy trading
reward, and our empirical study (Section~\ref{sec:signal}) independently rediscovers
the strategy literature's central lesson --- that predictive signal is scarce and is
best handled by risk management rather than by ever more flexible learning.

The objective an agent maximizes defines the behavior it learns; reward design is
therefore the central lever for risk. The canonical risk-adjusted choice, the
Differential Sharpe Ratio (DSR)~\cite{moody2001}, measures reward per unit of
\emph{variance}. Variance penalizes upside and downside symmetrically and so cannot
distinguish a strategy that occasionally produces large gains from one that
occasionally crashes. Yet crash risk --- negative skewness and heavy tails --- is
precisely what a risk-averse investor most wishes to avoid.

We propose replacing variance-based risk adjustment with the \emph{entropic certainty
equivalent}, the certainty equivalent of an exponential-utility investor and a
classical convex risk measure~\cite{follmer2002}, whose Taylor expansion weights every
cumulant of the return distribution through a single interpretable parameter $\theta$.
We turn this static measure into an online, differentiable per-step reward --- the
Differential Entropic Reward (DER) --- suitable for gradient-based RRL, and we situate
it within the risk-sensitive stochastic control tradition~\cite{bielecki1999,
flemingmce1995} and the modern risk-sensitive RL regret theory built on the same
exponential criterion~\cite{fei2020,fei2021,noorani2022}.

\paragraph{Contributions.}
\begin{enumerate}
\item \textbf{The DER reward and its theory} (Section~\ref{sec:der}): limiting cases,
a cumulant expansion, an exact telescoping identity, a KL-robust dual, and a
risk-sensitive HJB interpretation, each confirmed numerically
(Section~\ref{sec:verify}).
\item \textbf{A controlled experiment} isolating the moment-sensitivity that
distinguishes the DER from any Sharpe-based reward (Section~\ref{sec:controlled}).
\item \textbf{A cross-sectional backtest} on $48$ real equities establishing
statistically decisive drawdown dominance over the DSR (Section~\ref{sec:backtest}).
\item \textbf{Honest signal ablations} on a $470$-stock universe that delimit what the
reward can and cannot do, and that we argue is a more scientifically candid position
than prevailing SOTA reports (Section~\ref{sec:signal}).
\item \textbf{An architecture-agnostic head-to-head} showing the DER transfers
unchanged to actor--critic and PPO, where \emph{DER-PPO matches the risk control of
CVaR-PPO (CPPO) at roughly a quarter of the compute and without LLM signals}
(Section~\ref{sec:ppo}).
\item \textbf{A Fundamental-Law decomposition} that locates the DER precisely
($IR=IC\cdot\sqrt{BR}\cdot TC$): the reward contributes to the transfer/risk term, not
to forecasting skill, and a combined-signal engine wrapped by the DER is the honest
route to (modest) alpha (Section~\ref{sec:alpha}).
\end{enumerate}
All experiments use real S\&P~500 data and are fully reproducible
(Section~\ref{sec:repro}).

\section{Related Work}\label{sec:related}

\paragraph{RL for trading.} Direct policy search for trading originates with Moody and
Saffell's RRL and the Differential Sharpe Ratio~\cite{moody2001}, later extended with
deep representations~\cite{deng2017}, deep $Q$-learning agents (TDQN)~\cite{theate2021},
and end-to-end portfolio networks~\cite{jiang2017}. Task-specific agents address
execution~\cite{nevmyvaka2006} and market making~\cite{spooner2018}. The FinRL
ecosystem~\cite{finrl2020} standardized environments, data pipelines, and baselines,
and remains the dominant open-source framework. Our agent adopts the RRL formulation
of~\cite{moody2001}; our contribution is orthogonal to the architecture and concerns
the \emph{reward}.

\paragraph{Risk measures and risk-sensitive control.} Coherent and convex risk
measures~\cite{follmer2002} provide the axiomatic foundation; the entropic risk
measure --- the exponential-utility certainty equivalent we adopt --- is the canonical
convex measure with a Kullback--Leibler dual. In continuous time, maximizing the
entropic criterion is a risk-sensitive stochastic control
problem~\cite{bielecki1999,flemingmce1995}, whose HJB equation carries an
exposure-penalizing term governed by the risk-aversion parameter. Conditional
Value-at-Risk (CVaR)~\cite{rockafellar2000} is the competing risk functional in
applied RL; unlike the entropic measure it is not smooth in its quantile and, in RL,
typically requires sample sorting or an auxiliary distributional critic.

\paragraph{Risk-sensitive RL.} The exponential/entropic criterion underlies the
provably efficient risk-sensitive RL of Fei et al.~\cite{fei2020,fei2021}, the
exponential-criteria actor--critic of Noorani et al.~\cite{noorani2022}, and the
risk-sensitive $Q$-learning of Borkar~\cite{borkar2002}; CVaR appears in safety- and
distributionally-constrained RL. These works establish the \emph{statistical} theory
of the entropic objective (regret bounds, exploration). We complement them on the
\emph{reward-engineering} side: the DER renders the same exponential criterion as an
online, differentiable, telescoping per-step reward for direct-policy trading agents,
where the relevant baseline is the DSR rather than tabular value iteration.

\paragraph{Current SOTA and positioning.} The most recent risk-sensitive trading agents
infuse LLM-generated signals into CVaR-PPO (CPPO)~\cite{benhenda2025} and its
successors~\cite{dapo2025}. These report \emph{regime-dependent} outcomes --- one
variant dominating in bull markets, another in bear markets --- and explicitly
acknowledge persistent large drawdowns and limited interpretability as open problems.
Against this backdrop, the DER offers three differentiators: (i) it is a closed-form,
single-parameter reward requiring no LLM, news feed, auxiliary critic, or sample
sorting; (ii) it is grounded in a classical risk measure and the risk-sensitive HJB
equation rather than assembled heuristically; and (iii) we report its limits honestly,
demonstrating decisive drawdown control while declining to claim alpha that our
ablations show does not exist.

\section{Problem Formulation}\label{sec:formulation}

We consider a single risky asset with price $p_t$ and one-period return
$r_t=(p_t-p_{t-1})/p_{t-1}$. At each step the agent chooses a position
$F_t\in[-1,1]$ (long, short, or flat) via a policy $F_t=\tanh(w^\top x_t+b)$, where
$x_t$ is an observed state and $(w,b)$ are learned parameters. With proportional
transaction cost $\delta$, the realized trading return is
\begin{equation}
R_t \;=\; F_{t-1}\,r_t \;-\; \delta\,\lvert F_t-F_{t-1}\rvert .
\label{eq:Rt}
\end{equation}
The agent maximizes a cumulative reward $\sum_t \rho_t$ by gradient ascent on
$(w,b)$; the design question is the per-step reward $\rho_t$. The DSR
baseline~\cite{moody2001} maintains exponential moving averages $A_t,B_t$ of $R_t$
and $R_t^2$ and rewards the marginal change of the Sharpe ratio
$S_t=A_t/\sqrt{B_t-A_t^2}$, namely
$\DSR_t=(B_{t-1}\Delta A_t-\tfrac12A_{t-1}\Delta B_t)/(B_{t-1}-A_{t-1}^2)^{3/2}$
with $\Delta A_t=R_t-A_{t-1}$, $\Delta B_t=R_t^2-B_{t-1}$. We now construct an
alternative reward that is sensitive to all moments.

\section{The Differential Entropic Reward}\label{sec:der}

\subsection{Definition}

For an exponential-utility investor with risk aversion $\theta>0$, the certainty
equivalent of a random return $R$ is the entropic risk measure~\cite{follmer2002}
\begin{equation}
\CE_\theta(R) \;=\; -\frac{1}{\theta}\,\log \E\!\left[e^{-\theta R}\right].
\label{eq:ce}
\end{equation}

\begin{definition}[Differential Entropic Reward]
Let $M_t=(1-\eta)M_{t-1}+\eta\,e^{-\theta R_t}$ be an exponential moving average of
$e^{-\theta R_t}$ with rate $\eta\in(0,1)$. The Differential Entropic Reward is
\begin{equation}
\DER_t \;=\; -\frac{1}{\theta}\,\log\frac{M_t}{M_{t-1}} .
\label{eq:der}
\end{equation}
\end{definition}

The estimator $M_t$ is the online analogue of $\E[e^{-\theta R}]$ and $-\tfrac1\theta\log M_t$
the running certainty equivalent; $\DER_t$ rewards its one-step improvement.
Algorithm~\ref{alg:der} gives the full training loop.

\subsection{Properties}

\begin{theorem}[Limiting cases and moment expansion]\label{thm:expansion}
Let $R$ have finite cumulants $\kappa_1=\mu$, $\kappa_2=\sigma^2$, $\kappa_3$,
$\kappa_4$. Then
\begin{equation}
\CE_\theta(R)
=\mu-\frac{\theta}{2}\sigma^2+\frac{\theta^2}{6}\kappa_3-\frac{\theta^3}{24}\kappa_4+O(\theta^4).
\label{eq:cumexp}
\end{equation}
In particular $\CE_\theta(R)\to\mu$ as $\theta\to0$ (risk-neutral limit); to first
order $\CE_\theta(R)=\mu-\tfrac{\theta}{2}\sigma^2$, a mean--variance criterion.
\end{theorem}
\begin{proof}
By~\eqref{eq:ce}, $-\theta\,\CE_\theta=\log\E[e^{-\theta R}]=K_R(-\theta)$ is the
cumulant generating function evaluated at $-\theta$, with
$K_R(s)=\sum_{n\ge1}\kappa_n s^n/n!$ wherever the moment generating function is finite
in a neighborhood of $0$. Substituting $s=-\theta$ gives
$-\theta\,\CE_\theta=-\theta\kappa_1+\tfrac{\theta^2}{2}\kappa_2-\tfrac{\theta^3}{6}\kappa_3+\tfrac{\theta^4}{24}\kappa_4-\cdots$;
dividing by $-\theta$ yields~\eqref{eq:cumexp}. The limits follow by truncation and
continuity.
\end{proof}

The signs are economically meaningful: $\CE_\theta$ rewards mean and positive skewness
while penalizing variance and kurtosis, with weights scaling as increasing powers of
$\theta$. The Sharpe ratio depends only on $(\kappa_1,\kappa_2)$ and is by construction
blind to $\kappa_3,\kappa_4$.

\begin{proposition}[Exact telescoping]\label{prop:tele}
With $\CE_t\coloneq-\tfrac1\theta\log M_t$, the per-step rewards sum to the change in
the running certainty equivalent: $\sum_{t=1}^{T}\DER_t = \CE_T-\CE_0$.
\end{proposition}
\begin{proof}
By~\eqref{eq:der}, $\sum_{t=1}^{T}\DER_t
=-\tfrac1\theta\sum_{t=1}^T(\log M_t-\log M_{t-1})
=-\tfrac1\theta(\log M_T-\log M_0)=\CE_T-\CE_0$.
\end{proof}

Proposition~\ref{prop:tele} is the analogue of the DSR's defining property and ensures
that maximizing the sum of myopic per-step rewards maximizes the terminal
risk-adjusted objective, so the reward is consistent with the agent's long-run goal.

\begin{proposition}[Robust / worst-case dual]\label{prop:dual}
The entropic certainty equivalent admits the variational representation
\begin{equation}
\CE_\theta(R)=\inf_{Q\ll P}\Big\{\E_Q[R]+\tfrac{1}{\theta}\,\KL(Q\Vert P)\Big\},
\label{eq:dual}
\end{equation}
where $P$ is the reference measure and the infimum is over measures $Q$ absolutely
continuous w.r.t.\ $P$.
\end{proposition}
\begin{proof}[Proof sketch]
This is the Donsker--Varadhan / Gibbs variational identity, equivalently the dual
representation of the entropic convex risk measure~\cite{follmer2002}. The infimum is
attained at the tilted measure $dQ^\star/dP\propto e^{-\theta R}$, for which
$\E_{Q^\star}[R]+\tfrac1\theta\KL(Q^\star\Vert P)=-\tfrac1\theta\log\E_P[e^{-\theta R}]$.
\end{proof}

Equation~\eqref{eq:dual} gives the third reading of $\theta$: the agent optimizes
against the worst plausible measure $Q$, penalized by its Kullback--Leibler divergence
from the estimated $P$; larger $\theta$ admits more adversarial $Q$, i.e.\ greater
distrust of the model.

\subsection{Risk-sensitive HJB interpretation}

In continuous time, maximizing $\CE_\theta$ of terminal wealth is a risk-sensitive
stochastic control problem~\cite{bielecki1999,flemingmce1995}. For controlled state
$dX=b(X,u)\,dt+\sigma(X)\,dW$ and value function $V$, the risk-sensitive HJB equation is
\begin{equation}
\partial_t V+\sup_u\Big[b^\top\nabla V+\tfrac12\,\mathrm{tr}\!\big(\sigma\sigma^\top\nabla^2V\big)
-\tfrac{\theta}{2}\,\lVert\sigma^\top\nabla V\rVert^2\Big]=0,
\label{eq:hjb}
\end{equation}
whose final term penalizes exposure to volatility and vanishes as $\theta\to0$,
recovering the risk-neutral equation. The DER is the discrete-time, model-free shadow
of~\eqref{eq:hjb}: rather than solving the HJB with a known model, the agent ascends
the empirical $\sum_t\DER_t$. This grounding distinguishes the DER from ad-hoc
composite rewards.

\begin{figure}[t]
\centering
\fbox{\begin{minipage}{0.92\textwidth}
\small
\textbf{Algorithm 1: DER-RRL training}\\[2pt]
\textbf{Input:} returns $\{r_t\}$, features $\{x_t\}$, cost $\delta$, risk aversion
$\theta$, EMA rate $\eta$, learning rate $\alpha$, epochs $E$.\\
Initialize policy parameters $(w,b)$.\\
\textbf{for} epoch $=1,\dots,E$ \textbf{do}\\
\hspace*{1.5em} compute positions $F_t=\tanh(w^\top x_t+b)$ and returns
$R_t=F_{t-1}r_t-\delta|F_t-F_{t-1}|$ \eqref{eq:Rt};\\
\hspace*{1.5em} update $M_t=(1-\eta)M_{t-1}+\eta e^{-\theta R_t}$ and
$\DER_t=-\tfrac1\theta\log(M_t/M_{t-1})$ \eqref{eq:der};\\
\hspace*{1.5em} objective $J(w,b)=\sum_t \DER_t = \CE_T-\CE_0$
(Prop.~\ref{prop:tele});\\
\hspace*{1.5em} $(w,b)\leftarrow(w,b)+\alpha\,\nabla_{(w,b)}J$ \ (Adam).\\
\textbf{return} $(w,b)$.
\end{minipage}}
\caption{The DER replaces only the per-step reward; the RRL machinery is unchanged.
All maps ($\tanh$, $e^{(\cdot)}$, $\log$) are smooth, so $J$ is differentiable and
trainable by gradient ascent.}
\label{alg:der}
\end{figure}

\section{Numerical Verification of the Theory}\label{sec:verify}

Following a verify-before-claiming protocol, each theoretical statement is checked by
two independent computations (the exact entropic formula~\eqref{eq:ce} versus an
independent construction), seed $2025$. Table~\ref{tab:verify} summarizes; all eight
checks pass.

\begin{table}[h]
\centering
\small
\begin{tabular}{lll}
\toprule
Claim & Check & Result \\
\midrule
C1 risk-neutral limit & $\CE_{\theta\to0}\!\to\!\mu$ & agree to $10^{-5}$ \\
C2 first-order mean--variance & $\CE_\theta\!\approx\!\mu-\tfrac\theta2\sigma^2$ & agree to $5\times10^{-4}$ \\
C3 cumulant expansion (4 terms) & \eqref{eq:cumexp} vs exact & err $<6\times10^{-10}$ ($\theta{=}2$) \\
C4 telescoping & $\sum_t\DER_t=\CE_T-\CE_0$ & agree to $10^{-9}$ \\
C5 risk aversion & fair gamble $\pm0.05$, $\CE_{100}<0$ & $-0.043$, monotone in $\theta$ \\
C6 controlled example & equal Sharpe, $\CE$ differs & see Section~\ref{sec:controlled} \\
\bottomrule
\end{tabular}
\caption{Numerical verification of the DER theory (\texttt{verify\_der\_theory.py},
seed $2025$, $8/8$ pass).}
\label{tab:verify}
\end{table}

\section{Experimental Methodology}\label{sec:method}

\paragraph{Data.} Real daily OHLCV for S\&P~500 constituents, 2013-02-08 to 2018-02-07
($1259$ trading days), including the 2015--2016 correction. The single-asset and
cross-sectional backtests use $48$ and $470$ fully-covered constituents respectively.

\paragraph{Policy and features.} Unless noted, the state $x_t$ is the vector of the
most recent returns; the policy is the linear-$\tanh$ map of
Section~\ref{sec:formulation}. The cross-sectional study (Section~\ref{sec:signal})
uses per-day cross-sectionally standardized features (momentum, reversal,
volatility, distance-to-high) and market-neutral, gross-normalized weights.

\paragraph{Training and evaluation.} Parameters are trained by Adam on the
DER (or DSR) objective; transaction cost $\delta=5$\,bps; EMA rate $\eta$ as in the
RRL baseline. We use a $70/30$ chronological train/test split and report
\emph{only} out-of-sample test metrics, averaged over multiple random seeds.
Metrics are annualized Sharpe and Sortino ($\times\sqrt{252}$), terminal return, and
maximum drawdown (MDD). Paired comparisons use the Wilcoxon signed-rank test.

\section{Results}

\subsection{Controlled experiment: variance is not enough}\label{sec:controlled}

We construct two return streams with \emph{identical} mean and standard deviation ---
hence identical Sharpe ratio ($0.0667$) --- but opposite skewness: one with rare large
gains, one with rare large crashes (Figure~\ref{fig:motiv}). The DSR cannot prefer
one. The entropic criterion discriminates sharply: at $\theta=100$ the favorable
stream has $\CE=-0.0030$ versus $\CE=-0.0231$ for the crash-prone stream (claim C6).
The gap is entirely due to the third- and higher-order terms in~\eqref{eq:cumexp}.

\begin{figure}[t]
\centering
\includegraphics[width=0.86\textwidth]{motivating.png}
\caption{Two streams with identical mean, variance, and Sharpe ratio but opposite
skewness. The Sharpe ratio is identical by construction; the entropic certainty
equivalent separates them, and the gap widens with $\theta$.}
\label{fig:motiv}
\end{figure}

\subsection{Single-asset and regime behavior}\label{sec:single}

On individual equities, the risk-aversion parameter behaves as a monotone risk dial.
For a declining stock (GE), increasing $\theta$ monotonically tightens drawdown while
shrinking exposure (Table~\ref{tab:thetasweep}). Against the DSR on declining names,
the DER's protection is stark: on XOM the DSR returns $-36.6\%$ with a $-40.3\%$
drawdown, whereas the DER returns $-0.8\%$ with a $-1.0\%$ drawdown; on GE the DSR
drawdown is $-19.4\%$ versus $-2.2\%$ for the DER.

\begin{table}[h]
\centering
\small
\begin{tabular}{lrrrr}
\toprule
$\theta$ & 25 & 75 & 150 & 300 \\
\midrule
Max drawdown & $-6.4\%$ & $-2.2\%$ & $-1.1\%$ & $-0.6\%$ \\
\bottomrule
\end{tabular}
\caption{DER on GE (out-of-sample): $\theta$ is a monotone risk dial.}
\label{tab:thetasweep}
\end{table}

\subsection{Cross-sectional backtest}\label{sec:backtest}

We train single-asset RRL agents with the DSR and the DER ($\theta=75$) on $48$
S\&P~500 constituents, three seeds each, out-of-sample. Table~\ref{tab:backtest}
aggregates the test set; Figure~\ref{fig:backtest} shows the per-asset distributions.

\begin{table}[h]
\centering
\begin{tabular}{lrrr}
\toprule
Strategy & Mean return & Mean Sharpe & Mean MaxDD \\
\midrule
Buy \& Hold        & $+41.1\%$ & $0.95$ & $-17.1\%$ \\
RRL + DSR          & $+6.5\%$  & $0.05$ & $-23.2\%$ \\
RRL + DER ($\theta=75$) & $+0.6\%$ & $0.18$ & $\mathbf{-1.7\%}$ \\
\bottomrule
\end{tabular}
\caption{Out-of-sample aggregates across $48$ real S\&P~500 equities.}
\label{tab:backtest}
\end{table}

\begin{figure}[t]
\centering
\includegraphics[width=0.92\textwidth]{backtest_plots.png}
\caption{Left: distribution of maximum drawdown across $48$ equities --- DER (red)
concentrates near zero, DSR (blue) spreads to $-50\%$. Right: per-asset Sharpe,
DER vs DSR.}
\label{fig:backtest}
\end{figure}

The headline result concerns risk. The DER yields a smaller maximum drawdown than the
DSR for \emph{all $48$ assets} ($100\%$); no DER agent exceeds a $-20\%$ drawdown,
whereas $54\%$ of DSR agents do; the worst DSR drawdown is $-50.4\%$ versus $-8.3\%$
for the DER. A paired Wilcoxon signed-rank test on maximum drawdown gives
$p\approx7\times10^{-15}$. The Sharpe-ratio difference favors the DER in only $56\%$
of assets and is \emph{not} significant ($p\approx0.18$). The DER's advantage is thus
specifically and decisively in downside control --- consistent with its theory --- not
in risk-adjusted return per se. As expected in a bull sample, a passive long position
attains the highest raw Sharpe; no active single-asset strategy on a linear policy
beats it, a point we return to below.

\subsection{What the reward cannot do: honest signal ablations}\label{sec:signal}

A reward shapes \emph{how} an agent trades a given signal; it cannot create signal
absent from the data. We document this boundary on a market-neutral cross-sectional
universe of $470$ equities.

\paragraph{Technical indicators.} Adding RSI, MACD, Bollinger \%B and related features
helps on some assets (AAPL test Sharpe rises from $0.46$ to $1.35$) but not others; the
effect is asset- and period-dependent and does not survive aggregation. Such indicators
are deterministic functions of price and add no new information.

\paragraph{Volume / order-flow.} We constructed six causal volume signals
(volume-confirmed move, Amihud illiquidity, volume shock, OBV slope, price--volume
correlation, money-flow index). Their out-of-sample information coefficients are all
indistinguishable from zero ($|t|<2$; combined IC $=0.0001$), and four of the six are
mutually redundant (cross-sectional correlation $>0.88$). Volume is an input to
\emph{risk}, not to direction, at the daily horizon --- which is precisely where a
state-dependent $\theta$ could use it.

\paragraph{Signal learning overfits.} A fixed $12{-}1$ momentum factor earns an
out-of-sample Sharpe of $0.68$, uncorrelated with the market. Attempting to
\emph{learn} the cross-sectional score --- with an MLP (test Sharpe $0.08$) or an
$L_2$-regularized linear model (test Sharpe $-1.26$, with near-zero momentum loading)
--- robustly destroys this signal through overfitting; the fixed, theory-driven factor
outperforms every learned model. Adding the momentum factor as a market-neutral
overlay to a long market position improves the combined Sharpe modestly
($1.63\to1.65$) and reduces drawdown ($-7.6\%\to-4.9\%$), a diversification rather
than an alpha effect (Figure~\ref{fig:levers}).

\begin{figure}[t]
\centering
\includegraphics[width=0.86\textwidth]{all_levers_final.png}
\caption{Out-of-sample equity ($470$ equities). A fixed cross-sectional momentum
overlay added to the market improves risk-adjusted return and reduces drawdown;
\emph{learning} the signal (not shown) underperforms the fixed factor.}
\label{fig:levers}
\end{figure}

\paragraph{Implication.} These are boundaries, not failures. They establish that the
contribution of the DER is risk shaping on top of a fixed signal, and that performance
gains must come from new orthogonal information, not from greater model flexibility.

\section{Architecture-Agnostic Risk Control: Actor--Critic and PPO}\label{sec:ppo}

The DER is a reward, not an architecture; it should therefore transfer to richer
agents. We verify this and conduct a head-to-head against the CVaR objective used by
current state-of-the-art risk-sensitive agents~\cite{benhenda2025,rockafellar2000}.

\paragraph{Risk-sensitive critic.} In an actor--critic agent the critic learns a value
function. The risk-neutral Bellman equation $V(x_t)=\E[R_t+\gamma V(x_{t+1})]$ becomes,
under the exponential criterion, the risk-sensitive Bellman
equation~\cite{fei2021,borkar2002}
\begin{equation}
V_\theta(x_t)=-\tfrac1\theta\log\E\!\big[e^{-\theta(R_t+\gamma V_\theta(x_{t+1}))}\big],
\label{eq:rsbellman}
\end{equation}
which reduces to the risk-neutral form as $\theta\to0$. The cleanest and most robust
way to realize~\eqref{eq:rsbellman} in practice is to feed the DER as the per-step
reward: the agent then maximizes cumulative DER, and risk sensitivity enters through
the reward rather than through bespoke critic machinery.

\paragraph{Actor--critic transfer.} On the regime-switching environment with crashes
(Section~\ref{sec:single}), an advantage actor--critic (A2C) trained on the DER reward
controls drawdown far better than the same A2C trained on raw P\&L or the DSR: mean
out-of-sample MaxDD of $-3.8\%$ for DER versus $-29.5\%$ (risk-neutral) and $-24.2\%$
(DSR), across three seeds. The drawdown-control signature is thus a property of the
reward and survives the change of architecture.

\paragraph{PPO head-to-head: DER vs CVaR.} We then fix a single PPO core (clipped
surrogate, shared per-asset Gaussian policy, long-only exposure) and vary \emph{only}
the risk objective across three reward streams: raw P\&L (risk-neutral PPO), a
Rockafellar--Uryasev tail penalty (CPPO, i.e.\ CVaR), and the DER (DER-PPO). The
universe is a $55$-name Nasdaq-100 subset, out-of-sample, three seeds.
Table~\ref{tab:ppo} and Figure~\ref{fig:ppo} report the results.

\begin{table}[h]
\centering
\begin{tabular}{lrrrr}
\toprule
Strategy & Return & Sharpe & Sortino & MaxDD \\
\midrule
Equal-weight (B\&H)   & $+41.2\%$ & $2.13$ & $2.52$ & $-7.9\%$ \\
PPO (risk-neutral)    & $+18.9\%$ & $2.10$ & $2.49$ & $-4.0\%$ \\
CPPO (CVaR)           & $+18.2\%$ & $2.11$ & $2.49$ & $-3.8\%$ \\
DER-PPO ($\theta=50$) & $+18.7\%$ & $2.11$ & $2.49$ & $-3.9\%$ \\
\bottomrule
\end{tabular}
\caption{PPO head-to-head, out-of-sample (Nasdaq-100 subset). DER-PPO is statistically
indistinguishable from CPPO in risk-adjusted terms.}
\label{tab:ppo}
\end{table}

\begin{figure}[t]
\centering
\includegraphics[width=0.95\textwidth]{ppo_headtohead.png}
\caption{Left: out-of-sample equity --- the three PPO variants are nearly
indistinguishable and all reduce drawdown relative to the passive benchmark. Right:
training cost --- DER-PPO attains the same risk control as CPPO at roughly a quarter
of the compute, because the entropic reward is a closed-form exponential moving average
whereas CVaR requires per-step rolling-quantile estimation.}
\label{fig:ppo}
\end{figure}

Two findings stand out. First, \emph{DER-PPO matches CPPO}: their risk-adjusted
performance is identical to within seed noise (Sharpe $2.11$ versus $2.11$, MaxDD
$-3.9\%$ versus $-3.8\%$), yet the DER reward is closed-form and differentiable,
trains in roughly a quarter of the wall-clock time of the CVaR variant (no per-step
quantile sorting), and requires no LLM or news signals of the kind used in current
SOTA~\cite{benhenda2025}. Second, and honestly, on this benign, largely
trending test window the three risk objectives \emph{converge} --- there is little
tail risk to differentiate them; the DER's decisive separation from the risk-neutral
objective appears under stress, as established in Sections~\ref{sec:single}
and~\ref{sec:backtest}. A fully stressed head-to-head (a test period containing a
2020- or 2022-style crash) is the natural confirmation and is left to future work, as
it requires data beyond our cached window.

\section{Where Does Alpha Come From? A Fundamental-Law Decomposition}\label{sec:alpha}

The preceding sections show that the DER controls risk but does not, by itself,
generate excess return --- a property our ablations (Section~\ref{sec:signal}) trace
to the scarcity of predictive signal. To make this precise and constructive, we adopt
the Fundamental Law of Active Management~\cite{grinold1989,clarke2002}, which decomposes
the information ratio (IR, active return per unit of active risk) of any active strategy
as
\begin{equation}
IR \;=\; IC \cdot \sqrt{BR} \cdot TC,
\label{eq:funlaw}
\end{equation}
where $IC$ (information coefficient) is the cross-sectional correlation between a signal
and subsequent returns --- the forecasting skill; $BR$ (breadth) is the number of
independent bets per year; and $TC$ (transfer coefficient) is the efficiency with which
the signal is translated into positions under constraints and costs ($TC=1$ being
frictionless).

Equation~\eqref{eq:funlaw} locates the DER exactly: it contributes nothing to $IC$ ---
no reward can manufacture forecasting skill --- but it operates on $TC$ and on the risk
denominator, defending realized performance through drawdown control. Alpha must
therefore originate in $IC$ (a genuine signal) and breadth; the DER is the risk and
transfer layer placed on top.

We test this on the $470$-stock universe, out-of-sample, with monthly rebalancing. For
each candidate signal we measure the cross-sectional $IC$ and its significance, the
Fundamental-Law upper bound on IR (using $BR=N\times$ rebalances/yr, the optimistic
independent-bets assumption), and the realized IR of a long-short decile portfolio net
of $5$ bps. Table~\ref{tab:ic} reports the results.

\begin{table}[h]
\centering
\begin{tabular}{lrrrrr}
\toprule
Signal & $IC$ & $IC$ $t$-stat & pred.\ IR & real.\ IR & $TC$ \\
\midrule
12-1 momentum       & $0.024$ & $0.6$ & $1.82$ & $0.09$ & $0.05$ \\
6-1 momentum        & $0.038$ & $0.8$ & $2.83$ & $0.66$ & $0.24$ \\
1-month reversal    & $0.036$ & $0.9$ & $2.68$ & $0.41$ & $0.15$ \\
low volatility      & $0.004$ & $0.1$ & $0.27$ & $-0.75$ & --- \\
long trend (50/200) & $0.041$ & $0.8$ & $3.07$ & $0.45$ & $0.15$ \\
\midrule
\textbf{Combined}   & $0.041$ & $1.5$ & $3.04$ & $\mathbf{0.84}$ & $0.28$ \\
\bottomrule
\end{tabular}
\caption{Fundamental-Law decomposition (470-stock universe, out-of-sample, monthly
rebalancing). Predicted IR uses the optimistic independent-bets breadth; realized IR is
a long-short decile portfolio net of $5$ bps. Alpha is weak but positive and improves
with signal combination; the DER does not enter $IC$.}
\label{tab:ic}
\end{table}

Three lessons follow. First, individual signals carry only weak, statistically
insignificant skill ($IC\approx0.02$--$0.04$, all $|t|<2$), consistent with our
overfitting findings and with the intrinsic difficulty of return prediction. Second,
combining orthogonal signals raises both skill and realized performance (combined
realized IR $\approx0.84$, market-neutral, after costs) --- the breadth mechanism
of~\eqref{eq:funlaw}. Third, the gap between the optimistic predicted IR and the
realized IR is large ($TC\approx0.05$--$0.28$), reflecting the correlation of bets and
transaction costs; the independent-bets breadth is, as is well known, an upper
bound~\cite{clarke2002}.

The constructive reading is an architecture: a \emph{signal engine} that combines
weak-but-real orthogonal predictors to supply $IC$ and breadth, wrapped by the
\emph{DER as the risk layer} that defends the transfer coefficient and the drawdown.
This cleanly separates the two problems --- where return comes from (signal) and how
risk is controlled (the DER) --- and is the honest basis for any deployable system. We
caution that the present evidence is suggestive rather than conclusive: the $IC$s are
not individually significant over this single five-year window, and confirmation
requires a longer sample and strict walk-forward validation.

\subsection{Out-of-sample validation on a longer, stressed window
(2015--2024)}\label{sec:alpha-oos}

We now supply exactly that validation. Using free Yahoo Finance data --- which carries
its own caveats, discussed below --- we assemble a $150$-name large-capitalisation
universe over $2015$--$2024$, a window that, unlike the $2013$--$2018$ sample, contains
two genuine stress events (the March $2020$ COVID crash and the $2022$ bear market). We
introduce two \emph{new orthogonal signals} of the kind the Fundamental Law requires for
$IC$, and pass each through the identical gate of Section~\ref{sec:alpha}: (i)
post-earnings-announcement drift (PEAD)~\cite{bernard1989}, the standardized earnings
surprise entered with a strict $T{+}1$ delay after each announcement (no look-ahead),
and (ii) cross-sectional $12{-}1$ momentum, both in raw form and after demeaning
\emph{within} sector --- the sector-neutral construction that is the textbook remedy
for momentum crashes~\cite{daniel2016}. Risk is controlled by a state-dependent
$\theta_t$ driven by the VIX, the practical realisation of the DER risk layer.
Table~\ref{tab:ic-oos} and Figure~\ref{fig:gate-oos} report the gate;
Table~\ref{tab:strat-oos} and Figure~\ref{fig:equity-oos} the walk-forward strategies.

\begin{table}[h]
\centering
\begin{tabular}{lrrr}
\toprule
Signal (2015--2024, OOS) & $IC$ & $IC$ $t$-stat & real.\ IR \\
\midrule
PEAD earnings surprise   & $0.023$ & $1.1$ & $-0.60$ \\
12-1 momentum            & $0.040$ & $0.9$ & $-0.02$ \\
sector-neutral momentum  & $0.010$ & $0.3$ & $-0.58$ \\
\bottomrule
\end{tabular}
\caption{Fundamental-Law gate on the longer, stressed window. \emph{None} of the three
signals clears the $|t|>2$ significance bar. The cross-sectional momentum edge that was
weakly positive over $2013$--$2018$ --- where the \emph{combined} factor reached a
realized IR of $0.84$ (Table~\ref{tab:ic}) --- does not reappear, and
sector-neutralisation \emph{lowers} the raw momentum $IC$ further over this sample. The
honest reading is a null: no deployable directional skill is present in these free-data
signals here.}
\label{tab:ic-oos}
\end{table}

\begin{figure}[h]
\centering
\includegraphics[width=0.72\textwidth]{alpha_gate_oos.png}
\caption{The Fundamental-Law gate on the stressed window. Bars are the absolute
out-of-sample $IC$ $t$-statistic of each signal; none clears the $|t|>2$ threshold
(dashed line), so all are rejected (red). Information coefficients are annotated above
each bar.}
\label{fig:gate-oos}
\end{figure}

\begin{table}[h]
\centering
\begin{tabular}{lrrr}
\toprule
Strategy (2015--2024, OOS) & Return & Sharpe & MaxDD \\
\midrule
Market (long-only)            & $+30.5\%$ & $0.70$ & $-17.6\%$ \\
12-1 momentum                 & $-21.0\%$ & $-0.76$ & $-25.6\%$ \\
Sector-neutral momentum       & $-27.7\%$ & $-1.71$ & $-29.2\%$ \\
\quad + DER VIX-$\theta_t$ risk layer & $-23.3\%$ & $-1.50$ & $\mathbf{-24.8\%}$ \\
Market + market-neutral overlay & $+6.6\%$ & $0.27$ & $\mathbf{-14.3\%}$ \\
\bottomrule
\end{tabular}
\caption{Walk-forward strategies on the stressed window. The cross-sectional momentum book
\emph{loses} against a passive long position in this mega-cap-led regime --- a direct
out-of-sample refutation of the durability of the $2013$--$2018$ momentum edge. The DER
risk layer does what it is designed to do and nothing it is not: VIX-driven $\theta_t$
tightens the momentum book's drawdown ($-24.8\%$ vs.\ $-29.2\%$), and adding the
market-neutral overlay to the market reduces drawdown ($-14.3\%$ vs.\ $-17.6\%$). On the
worst rolling $60$-day sub-window the risk-layered book draws down $-9.0\%$ against the
market's $-16.0\%$.}
\label{tab:strat-oos}
\end{table}

\begin{figure}[h]
\centering
\includegraphics[width=0.85\textwidth]{sector_mom_oos.png}
\caption{Out-of-sample equity on the stressed window ($2015$--$2024$, including the
$2020$ and $2022$ stress events). The cross-sectional momentum book underperforms a
passive market position; the DER VIX-$\theta_t$ risk layer and the market-neutral overlay
deliver their effect through drawdown control rather than return.}
\label{fig:equity-oos}
\end{figure}

Two conclusions follow, and we state them without softening. First, the modest momentum
alpha documented over $2013$--$2018$ does \emph{not} survive into $2015$--$2024$: the
$IC$ is insignificant and the long--short book underperforms a passive position. This is
consistent with the well-documented post-$2015$ decay and crash-proneness of cross-sectional
momentum~\cite{daniel2016} and vindicates the caution of the preceding paragraph rather
than overturning it. Second, and in the same experiment, the DER risk layer continues to
deliver its one demonstrable effect --- drawdown control --- including on the rolling
stressed sub-window, exactly as the single-asset and PPO results predicted under stress.
The separation the Fundamental Law makes precise is therefore reinforced, not weakened, by
the longer sample: \emph{return is scarce and regime-dependent; risk control is robust and
transferable.}

\paragraph{Caveats of the free-data validation.} This validation inherits three
limitations we name explicitly rather than bury. The universe is built from
\emph{current} constituents, so it carries survivorship bias; the earnings consensus used
for the surprise is the vendor's latest figure rather than a strict point-in-time
snapshot, which can only \emph{flatter} PEAD --- yet PEAD still fails the gate; and the
breadth in $IR=IC\sqrt{BR}\,TC$ remains the optimistic independent-bets bound. Each of
these biases works \emph{in favour} of finding alpha, which makes the null result more,
not less, credible.

\section{Discussion}\label{sec:discussion}

The DER occupies a specific and, we argue, well-motivated point in design space.
Relative to the DSR~\cite{moody2001} it adds sensitivity to skewness and kurtosis at no
architectural cost and with a single interpretable parameter. Relative to
CVaR-based agents~\cite{rockafellar2000,benhenda2025} it is smooth, differentiable, and
critic-free, and inherits a closed-form dual~\eqref{eq:dual} and HJB
interpretation~\eqref{eq:hjb}. Relative to the risk-sensitive RL
theory~\cite{fei2020,fei2021,noorani2022}, which studies the same exponential criterion
in tabular or function-approximation regret settings, we contribute its
\emph{reward-engineering} realization for direct-policy trading. The empirical
signature --- decisive, statistically significant drawdown control without a claim of
alpha --- is, we contend, a more scientifically defensible report than regime-dependent
return claims, and is directly useful for capital that must be protected.

\section{Limitations and Future Work: Toward Risk-Managed Strategies}\label{sec:future}

\paragraph{Limitations.} The study is confined to daily equity data over a single
five-year window with linear or shallow policies. The directional-signal conclusions
may differ at intraday horizons, where genuine microstructure information
exists~\cite{nevmyvaka2006,spooner2018}, or with fundamental and cross-asset features
not accessible in our setting. Crucially, our ablations (Section~\ref{sec:signal})
establish that performance gains require new \emph{information}, not a new reward; we
therefore frame the following directions as deliberate future work rather than claims.

\paragraph{A research-grounded roadmap.} The strategy literature suggests a concrete,
testable programme in which the DER plays its proven role --- risk shaping --- on top
of documented signals.
\begin{enumerate}
\item \textbf{Orthogonal signal ensemble.} Rather than learning a single signal (which
we showed overfits), combine the few signals known to carry out-of-sample information:
cross-sectional momentum~\cite{jegadeesh1993} and time-series
momentum~\cite{moskowitz2012}, which is documented to ``perform best during extreme
markets.'' Each is used with fixed, theory-driven sign, and the ensemble is weighted
by \emph{risk parity}~\cite{asness2012} rather than by a learned mapping, directly
addressing the leverage-aversion and diversification concerns that motivate that
construction.
\item \textbf{State-dependent $\theta$ as principled volatility management.} Momentum's
characteristic failure mode is the \emph{momentum crash}~\cite{daniel2016}, and the
established remedy is to scale exposure by ex-ante inverse volatility. This is exactly
what a regime-aware, state-dependent risk aversion $\theta_t$ implements within the DER:
raising $\theta_t$ when volatility (or the directionally-uninformative but
risk-informative volume signals of Section~\ref{sec:signal}) is high de-risks the book
before turbulence. Volatility-managed portfolios~\cite{moreira2017} report higher
Sharpe ratios precisely because volatility is more forecastable than returns ---
giving the DER risk layer a documented mechanism through which to add value.
\item \textbf{Honest out-of-sample caveat.} We will test these against the known
caution that volatility management can \emph{harm} real-time performance and is most
reliable when combined with, rather than replacing, the unmanaged
strategy~\cite{moreira2017}; the evaluation must therefore be strictly walk-forward and
report combined as well as stand-alone variants.
\item \textbf{Methodological neighbours.} The Sharpe-optimised deep momentum networks
of Lim, Zohren and Roberts~\cite{lim2019}, which learn position sizing within a
volatility-scaling framework and add a turnover regulariser, are the closest
methodological reference point; replacing their Sharpe objective with the DER, and
their LSTM with our entropic per-step reward, is a natural head-to-head experiment.
\end{enumerate}
This programme keeps the contribution of the present paper intact --- the DER as a
principled, single-parameter risk layer --- while pursuing performance through the only
avenue our evidence supports: combining genuine, orthogonal signals under disciplined,
volatility-aware risk control.

\section{Conclusion}\label{sec:conclusion}

We introduced the Differential Entropic Reward, an online, differentiable, telescoping
per-step reward that renders the entropic certainty equivalent --- and hence
sensitivity to all return cumulants through a single parameter $\theta$ --- usable by
gradient-based RL trading agents. We proved its limiting cases, cumulant expansion,
telescoping identity, and robust dual, verified each numerically, and showed
empirically on real equities that it controls drawdown decisively and significantly
relative to the Differential Sharpe Ratio, while honestly delimiting that it does not
manufacture predictive signal. The DER is a principled, interpretable,
single-parameter mechanism for the one thing it demonstrably delivers: control of
downside risk.

\section*{Reproducibility}\label{sec:repro}
\addcontentsline{toc}{section}{Reproducibility}
\texttt{verify\_der\_theory.py} checks claims C1--C6 two ways (seed $2025$, $8/8$ pass).
\texttt{backtest.py} regenerates Table~\ref{tab:backtest}, Figure~\ref{fig:backtest},
and the Wilcoxon tests; \texttt{harness.py} runs any single-asset configuration and
Table~\ref{tab:thetasweep}; \texttt{volume\_signals.py} reproduces the
information-coefficient and redundancy analysis; \texttt{all\_levers\_v3.py} reproduces
the cross-sectional results and Figure~\ref{fig:levers}; \texttt{motivating\_example.py}
reproduces Figure~\ref{fig:motiv}. \texttt{actor\_critic\_der.py} reproduces the
actor--critic transfer result, and \texttt{ppo\_headtohead.py} with
\texttt{ppo\_figure.py} reproduce Table~\ref{tab:ppo} and Figure~\ref{fig:ppo}.
\texttt{ic\_analysis.py} reproduces the Fundamental-Law decomposition
(Table~\ref{tab:ic}). Data are real S\&P~500 OHLCV (2013--2018), auto-downloaded from
a public mirror.

\begin{thebibliography}{99}
\small
\bibitem{moody2001} J.~Moody and M.~Saffell. Learning to trade via direct
reinforcement. \emph{IEEE Transactions on Neural Networks}, 12(4):875--889, 2001.
\bibitem{deng2017} Y.~Deng, F.~Bao, Y.~Kong, Z.~Ren, and Q.~Dai. Deep direct
reinforcement learning for financial signal representation and trading. \emph{IEEE
Transactions on Neural Networks and Learning Systems}, 28(3):653--664, 2017.
\bibitem{theate2021} T.~Th\'eate and D.~Ernst. An application of deep reinforcement
learning to algorithmic trading. \emph{Expert Systems with Applications}, 173:114632,
2021.
\bibitem{jiang2017} Z.~Jiang, D.~Xu, and J.~Liang. A deep reinforcement learning
framework for the financial portfolio management problem. arXiv:1706.10059, 2017.
\bibitem{nevmyvaka2006} Y.~Nevmyvaka, Y.~Feng, and M.~Kearns. Reinforcement learning
for optimized trade execution. In \emph{Proc.\ 23rd Int.\ Conf.\ on Machine Learning
(ICML)}, 2006.
\bibitem{spooner2018} T.~Spooner, J.~Fearnley, R.~Savani, and A.~Koukorinis. Market
making via reinforcement learning. arXiv:1804.04216, 2018.
\bibitem{hambly2023} B.~Hambly, R.~Xu, and H.~Yang. Recent advances in reinforcement
learning in finance. \emph{Mathematical Finance}, 33(3):437--503, 2023.
\bibitem{finrl2020} X.-Y.~Liu, H.~Yang, Q.~Chen, R.~Zhang, L.~Yang, B.~Xiao, and
C.~D.~Wang. FinRL: A deep reinforcement learning library for automated stock trading
in quantitative finance. \emph{Deep RL Workshop, NeurIPS}, 2020.
\bibitem{benhenda2025} M.~Benhenda. FinRL-DeepSeek: LLM-infused risk-sensitive
reinforcement learning for trading agents. arXiv:2502.07393, 2025.
\bibitem{dapo2025} R.~Zha et al. A new DAPO algorithm for stock trading.
arXiv:2505.06408, 2025.
\bibitem{follmer2002} H.~F\"ollmer and A.~Schied. Convex measures of risk and trading
constraints. \emph{Finance and Stochastics}, 6(4):429--447, 2002.
doi:10.1007/s007800200072.
\bibitem{bielecki1999} T.~R.~Bielecki and S.~R.~Pliska. Risk-sensitive dynamic asset
management. \emph{Applied Mathematics \& Optimization}, 39:337--360, 1999.
doi:10.1007/s002459900110.
\bibitem{flemingmce1995} W.~H.~Fleming and W.~M.~McEneaney. Risk-sensitive control on
an infinite time horizon. \emph{SIAM Journal on Control and Optimization},
33(6):1881--1915, 1995.
\bibitem{rockafellar2000} R.~T.~Rockafellar and S.~Uryasev. Optimization of conditional
value-at-risk. \emph{Journal of Risk}, 3:21--41, 2000. doi:10.21314/JOR.2000.038.
\bibitem{fei2020} Y.~Fei, Z.~Yang, Y.~Chen, Z.~Wang, and Q.~Xie. Risk-sensitive
reinforcement learning: near-optimal risk-sample tradeoff in regret. In
\emph{Advances in Neural Information Processing Systems (NeurIPS)}, 2020.
\bibitem{fei2021} Y.~Fei, Z.~Yang, Y.~Chen, and Z.~Wang. Exponential Bellman equation
and improved regret bounds for risk-sensitive reinforcement learning. In
\emph{Advances in Neural Information Processing Systems (NeurIPS)}, 2021.
\bibitem{noorani2022} E.~Noorani, C.~N.~Mavridis, and J.~S.~Baras. Risk-sensitive
reinforcement learning with exponential criteria. arXiv:2212.09010, 2022.
\bibitem{borkar2002} V.~S.~Borkar. Q-learning for risk-sensitive control.
\emph{Mathematics of Operations Research}, 27(2):294--311, 2002.
\bibitem{jegadeesh1993} N.~Jegadeesh and S.~Titman. Returns to buying winners and
selling losers: implications for stock market efficiency. \emph{Journal of Finance},
48(1):65--91, 1993. doi:10.1111/j.1540-6261.1993.tb04702.x.
\bibitem{bernard1989} V.~L.~Bernard and J.~K.~Thomas. Post-earnings-announcement drift:
delayed price response or risk premium? \emph{Journal of Accounting Research},
27:1--36, 1989. doi:10.2307/2491062.
\bibitem{moskowitz2012} T.~J.~Moskowitz, Y.~H.~Ooi, and L.~H.~Pedersen. Time series
momentum. \emph{Journal of Financial Economics}, 104(2):228--250, 2012.
\bibitem{daniel2016} K.~Daniel and T.~J.~Moskowitz. Momentum crashes. \emph{Journal of
Financial Economics}, 122(2):221--247, 2016.
\bibitem{moreira2017} A.~Moreira and T.~Muir. Volatility-managed portfolios.
\emph{Journal of Finance}, 72(4):1611--1644, 2017. doi:10.1111/jofi.12513.
\bibitem{asness2012} C.~S.~Asness, A.~Frazzini, and L.~H.~Pedersen. Leverage aversion
and risk parity. \emph{Financial Analysts Journal}, 68(1):47--59, 2012.
\bibitem{lim2019} B.~Lim, S.~Zohren, and S.~Roberts. Enhancing time series momentum
strategies using deep neural networks. \emph{Journal of Financial Data Science},
1(4):19--38, 2019. arXiv:1904.04912.
\bibitem{grinold1989} R.~C.~Grinold. The fundamental law of active management.
\emph{Journal of Portfolio Management}, 15(3):30--37, 1989.
\bibitem{clarke2002} R.~Clarke, H.~de Silva, and S.~Thorley. Portfolio constraints and
the fundamental law of active management. \emph{Financial Analysts Journal},
58(5):48--66, 2002.
\end{thebibliography}

\end{document}
